\section{Arithmetic fallout: a second-moment recurrence}\label{sec:fiber-spectrum}

The threshold and conjugacy theorems do not require a detailed study of the entire fiber spectrum, but one unconditional statistic is useful later and illustrates the remaining local arithmetic: the second moment
\[
\Sigma_m:=\sum_{x\in X_m}d_m(x)^2,
\]
equivalently the collision probability of the visible law under uniform raw input.

\subsection{Second moment and collision probability}\label{sec:fiber-collision}

\begin{theorem}[Second-moment linear recurrence]\label{prop:second-moment-recurrence}
For all $m\ge 4$,
\begin{equation}\label{eq:sigma-recurrence}
\Sigma_m=2\,\Sigma_{m-1}+2\,\Sigma_{m-2}-2\,\Sigma_{m-3},
\end{equation}
with $\Sigma_1=2$, $\Sigma_2=6$, $\Sigma_3=14$.
\end{theorem}

\begin{proof}
The second moment counts raw pairs sharing a residue modulo $F_{m+2}$. Write $\epsilon_k:=\omega_k-\omega_k'\in\{-1,0,1\}$ with multiplicity weights $w(\pm1)=1$, $w(0)=2$, and define
\[
f(n,t):=\text{weighted count of }\epsilon\in\{-1,0,1\}^n\text{ with }\sum_{k=1}^n \epsilon_kF_{k+1}=t.
\]
Since
\[
\Bigl|\sum_{k=1}^m \epsilon_kF_{k+1}\Bigr|\le \sum_{k=1}^mF_{k+1}=F_{m+3}-2<2F_{m+2},
\]
only the values $0,\pm F_{m+2}$ contribute, so by symmetry
\begin{equation}\label{eq:sigma-split}
\Sigma_m=f(m,0)+2f(m,F_{m+2}).
\end{equation}

Conditioning on the last digit gives
\[
f(m,F_{m+2})=2f(m-1,F_{m+2})+f(m-1,F_m)+f(m-1,F_{m+3}),
\]
and the first and third terms vanish by the same range bound. Hence
\[
f(m,F_{m+2})=f(m-1,F_m).
\]
Conditioning once more at level $m-1$ yields
\[
f(m-1,F_m)=2f(m-2,F_m)+f(m-2,0),
\]
because the term at $2F_m$ is out of range for $m\ge4$. Using \eqref{eq:sigma-split} at window $m-2$, this is exactly $\Sigma_{m-2}$.

Likewise,
\[
f(m,0)=2f(m-1,0)+2f(m-1,F_{m+1})=f(m-1,0)+\Sigma_{m-1}.
\]
Writing $a_m:=f(m,0)$, one has $a_m=\Sigma_m-2\Sigma_{m-2}$ and $a_{m-1}=\Sigma_{m-1}-2\Sigma_{m-3}$. Eliminating $a_m$ gives \eqref{eq:sigma-recurrence}. The initial values are obtained by direct enumeration.
\end{proof}

\begin{proposition}[Uniform visible law and collision probability]\label{prop:uniform-collision}
Let $\mu_m^{\mathrm{unif}}$ be the uniform distribution on $\{0,1\}^m$, and let $\pi_m^{\mathrm{unif}}:=(\Fold_m)_*\mu_m^{\mathrm{unif}}$ be the induced distribution on $X_m$. Then for every $x\in X_m$,
\[
\pi_m^{\mathrm{unif}}(x)=\frac{d_m(x)}{2^m},
\]
and therefore
\begin{equation}\label{eq:col-unif-second-moment}
\Col_{\mathrm{unif}}(m):=\sum_{x\in X_m}\pi_m^{\mathrm{unif}}(x)^2
=\frac{1}{2^{2m}}\sum_{x\in X_m} d_m(x)^2.
\end{equation}
\end{proposition}

\begin{proof}
For $x\in X_m$, the fiber $\mathcal{F}_m(x)$ has cardinality $d_m(x)$ by definition. Since $\mu_m^{\mathrm{unif}}$ assigns mass $2^{-m}$ to each raw word,
\[
\pi_m^{\mathrm{unif}}(x)=\mu_m^{\mathrm{unif}}(\mathcal{F}_m(x))=\frac{d_m(x)}{2^m}.
\]
Substituting into $\Col_{\mathrm{unif}}(m)=\sum_x \pi_m^{\mathrm{unif}}(x)^2$ gives \eqref{eq:col-unif-second-moment}.
\end{proof}

